\chapter{Quello che il tempo cancella}
\label{cap:tafonomia}

I due capitoli precedenti hanno mostrato la distribuzione attuale dell'arte rupestre nel Veneto e le ragioni per sospettare che quella distribuzione rispecchi in parte la storia della ricerca. Ma c'è una terza variabile, distinta dal research bias e dalla reale frequentazione antica, che condiziona in modo indipendente ciò che oggi è possibile trovare: la conservazione nel tempo delle tracce stesse. Prima di chiedersi dove cercare, vale la pena chiedersi cosa sopravvive, e per quanto, e perché.

\section{Due storie di conservazione molto diverse}

Il confronto tra il deserto di Atacama e le Alpi venete non è solo un confronto tra territori più o meno esplorati. È anche, e prima ancora, un confronto tra due ambienti di conservazione radicalmente diversi.

Il deserto di Atacama è uno degli ambienti più aridi del pianeta. In larghe fasce della regione le precipitazioni annue sono inferiori a un millimetro; in alcune stazioni meteorologiche non si registra pioggia per decenni consecutivi. In queste condizioni, i principali agenti di degrado di un'incisione o di una pittura rupestre, che altrove operano senza sosta, qui sono quasi inattivi: i cicli di gelo e disgelo sono rari o assenti alle quote rilevanti; la vegetazione che colonizza le superfici rocciose è quasi inesistente; i licheni, che in ambienti umidi avanzano di qualche millimetro all'anno, qui crescono a velocità irrisorie. Una pittura rupestre in questo contesto non ha molti nemici. Può invecchiare senza essere cancellata.

Le Alpi venete, in confronto, sono un ambiente di degrado attivo. Le precipitazioni annue superano i 1.000 millimetri su gran parte del territorio montano; i cicli di gelo e disgelo sono frequenti in quota per diversi mesi all'anno; la vegetazione ricolonizza rapidamente ogni superficie esposta, compresi i licheni crostosi che aderiscono alla roccia e ne accelerano la disgregazione chimica; le pareti rocciose non protette sono soggette a stillicidio, scivolamento di depositi superficiali, distacco di frammenti per crioclastismo. Un'incisione poco profonda in questo contesto ha la vita breve, anche in assenza di qualsiasi intervento umano.

\section{L'asimmetria temporale}

C'è però un'asimmetria ancora più fondamentale tra i due confronti, che è facile trascurare se ci si concentra troppo sul confronto geografico: i due fenomeni non si collocano nello stesso periodo.

Le pitture e i petroglifi del deserto di Atacama, nella grande maggioranza dei casi studiati, risalgono ai millenni compresi tra 3.000 e 500 anni prima dell'anno zero. Le prime corrispondono all'incirca all'apogeo delle culture regionali pre-incaiche; le più recenti si sovrappongono al periodo di espansione dell'impero inca e ai primi secoli del contatto europeo. Anche i siti più antichi attribuiti alla tradizione rupestre atacamena non superano in genere i 4.000--5.000 anni prima dell'anno zero. E, aspetto cruciale, nell'intero deserto di Atacama non sono ancora state identificate con certezza pitture o incisioni databili al periodo in cui le Alpi venete ci interessano di più: il Tardoglaciale e il Mesolitico, tra 17.000 e 7.500 anni prima dell'anno zero.

Questo non significa che il territorio fosse disabitato nel Tardoglaciale. Le miniere di ocra nell'area di Taltal, come discusso nel primo capitolo, attestano una presenza umana che risale probabilmente a 12.000--10.000 anni prima dell'anno zero. Ma quella presenza, ammesso che abbia lasciato segni sulle rocce, non ha lasciato segni che oggi sappiamo riconoscere o attribuire con sicurezza a quel periodo. I motivi possono essere molti: quei segni esistevano e sono stati cancellati dal tempo; esistevano e sono stati sovrapposti da segni successivi; non sono stati ancora cercati con metodi adeguati; oppure, più semplicemente, la pratica di incidere e dipingere la roccia nel modo che oggi riconosciamo non era ancora diffusa in quella fase.

In altri termini, il confronto diretto tra la ricchezza di arte rupestre del deserto di Atacama e la scarsità di quella veneta non è del tutto simmetrico nemmeno sul piano temporale: la ricchezza atacamena si riferisce a un arco di 3.000--5.000 anni relativamente recenti e in un ambiente di conservazione eccezionale; la scarsità veneta si riferisce a un arco cronologico molto più lungo e in un ambiente di conservazione sfavorevole. Non sono la stessa cosa, e mescolarle in un unico confronto rischia di produrre un'impressione più netta di quanto i dati giustifichino.

\section{Tre modi di non trovare un sito}

Questa asimmetria porta a distinguere tre situazioni che in superficie sembrano identiche (il sito non c'è), ma che hanno cause e implicazioni molto diverse.

La prima è l'assenza reale: in questo luogo non è mai stato inciso o dipinto nulla di intenzionale, o comunque nulla che oggi rientrerebbe nella categoria di arte rupestre. Può dipendere da ragioni culturali (le comunità che frequentavano questo territorio non praticavano questa forma di espressione), da ragioni pratiche (le rocce disponibili non erano adatte alla tecnica usata) o semplicemente dal caso.

La seconda è l'assenza apparente: il segno c'era, ma è stato cancellato o reso illeggibile dai processi di degrado, dall'erosione, dalla copertura vegetale, dal lichenamento, da crolli o da riutilizzi successivi. In questo caso l'assenza non è nel passato, ma nella catena tra il passato e il presente.

La terza è l'assenza non ancora osservata: il segno c'è, ma non è ancora stato trovato, o perché nessuno ha cercato in modo sistematico in quel luogo, o perché le condizioni di visibilità al momento della ricerca erano sfavorevoli, o perché chi ha cercato usava metodi non adatti a quel tipo di traccia.

Il research bias discusso nel secondo capitolo di questo libro riguarda principalmente la terza categoria. Ma il problema della tafonomia rupestre, cioè del degrado e della cancellazione delle tracce nel tempo, riguarda la seconda. Le due categorie non si escludono: un'area può essere sia poco indagata (terza categoria) sia fisicamente sfavorevole alla conservazione (seconda categoria). Confonderle porta però a conclusioni diverse: nella seconda, anche la ricerca più intensa non restituirebbe i siti perduti, perché non ci sono più; nella terza, la ricerca sistematica può ancora cambiare il quadro.

\section{I processi di degrado specifici per l'arte rupestre alpina}

Per l'area alpina, i principali processi che minacciano la leggibilità di un'incisione o di una pittura rupestre all'aperto sono documentati, anche se raramente quantificati in modo sistematico.

Il crioclastismo, già discusso nel terzo capitolo a proposito della grotta di Fumane, è il processo più generale: l'acqua penetra nelle microfessure della roccia, ghiaccia, si espande, e stacca frammenti. Sulle pareti verticali esposte, questo processo può rendere illeggibile un'incisione in poche centinaia di anni, specialmente se poco profonda. Un segno inciso a quattro millimetri di profondità resiste molto meno di un segnato a dodici; ma anche i segni profondi, se la parete si frattura lungo il piano dell'incisione, spariscono con il distacco dello strato superiore.

La lichenizzazione è un secondo processo difficile da contrastare. I licheni crostosi, in particolare le specie del genere \textit{Rhizocarpon} comuni sulle rocce calcaree alpine, colonizzano le superfici a velocità lenta ma continua, crescendo all'interno degli strati superficiali della roccia e disgregandola chimicamente attraverso la produzione di acidi organici. Un pannello ricco di licheni diventa progressivamente illeggibile non perché sia coperto da uno strato opaco, ma perché la superficie stessa della roccia, con le incisioni su di essa, viene consumata e polverizzata.

L'erosione idrica agisce in modo diverso ma convergente: le precipitazioni abbondanti sciolguono la roccia calcarea, levigando le superfici e riducendo le differenze di rilievo che rendono leggibili le incisioni. Sulle superfici sub-orizzontali esposte all'acqua corrente il processo è particolarmente rapido; sulle pareti verticali è più lento, ma agisce comunque nel tempo.

Infine, la copertura da parte di sedimenti, detrito vegetale, muschi e suolo organico può rendere invisibile un pannello anche senza cancellarlo: se il piano di campagna si alza di pochi decimetri nel corso dei secoli, porzioni di parete prima visibili diventano sepolte, e i segni che ospitavano rimangono fisicamente intatti ma invisibili a chi cammina sulla superficie attuale.

\section{Cosa può ancora essere trovato, e cosa no}

Questa rassegna non intende concludere che la ricerca sia inutile. Intende chiarire, per onestà metodologica, che la ricerca può cambiare il quadro nella misura in cui le tracce sono fisicamente sopravvissute e fisicamente accessibili.

Per le fasi più antiche, quelle mesolitiche e tardoglaciali che costituiscono il periodo di maggiore interesse di questo libro, il problema della conservazione è più acuto che per le fasi più recenti. Un'ipotetica incisione risalente a 10.000 anni prima dell'anno zero ha avuto dieci millenni di cicli di gelo e disgelo, di lichenazione, di erosione idrica e di variazioni del piano di campagna per diventare illeggibile. La probabilità che sopravviva in forma leggibile è bassa, anche se non nulla: esistono casi, in altri contesti alpini, di incisioni attribuite al Mesolitico in base al contesto stratigrafico, non solo allo stile.

Per le fasi più recenti, dall'età del bronzo in poi, la finestra di sopravvivenza è più ampia. I siti noti del Garda e della Val d'Assa rientrano in questa categoria, e mostrano che la sopravvivenza, almeno parziale, è possibile.

Questa distinzione temporale ha una conseguenza pratica per il tipo di ricerca da fare: cercare arte rupestre mesolitica nelle Alpi venete richiede non solo un metodo di prospezione sistematico, come quello di Kompatscher e Hrozny Kompatscher discusso nella terza parte di questo libro, ma anche metodi di rilevamento adattati a tracce parzialmente degradate o parzialmente coperte, come l'imaging multispettrale, la fotogrammetria a luce radente o la rimozione controllata di sedimenti superficiali. E richiede di accettare in anticipo che una parte delle tracce, forse la maggior parte, non sarà più trovabile, indipendentemente dall'intensità della ricerca.

\section{Un'ipotesi alternativa che non va ignorata}

Fin qui questo capitolo ha trattato la tafonomia come il principale fattore che spiega l'assenza di arte rupestre mesolitica nelle Alpi venete: i segni c'erano, ma sono stati cancellati o sono irraggiungibili. È un'ipotesi ragionevole, ma non è l'unica possibile, e sarebbe intellettualmente disonesto non enunciarla chiaramente.

L'ipotesi alternativa è che i cacciatori mesolitici delle Alpi venete semplicemente non producessero arte rupestre all'aperto su larga scala, o che non la producessero affatto. Non perché fossero culturalmente inferiori o privi di capacità simboliche: Fumane dimostra il contrario per l'arte parietale in grotta; la sepoltura di Mondeval, con il suo corredo di ocre e conchiglie, dimostra pratiche rituali elaborate per il Mesolitico. Ma l'arte parietale in grotta e l'arte rupestre all'aperto sono pratiche diverse, che non si implicano necessariamente a vicenda.

In molte regioni d'Europa il Mesolitico è quasi del tutto privo di arte rupestre all'aperto, e questo non è spiegato solo dal degrado differenziale: alcune delle stesse rocce che ospiteranno incisioni dell'età del Bronzo erano già esposte e accessibili durante il Mesolitico. La ragione più probabile è che la tradizione di marcare le rocce all'aperto in modo visibile e permanente sia associata, nella maggior parte dei contesti europei, a società più sedentarie, che hanno interesse a segnalare in modo duraturo il proprio rapporto con un territorio specifico. Cacciatori altamente mobili possono usare altre forme di comunicazione simbolica, materiali deperibili (pitture su corteccia, incisioni su legno, segnali al suolo) che non lasciano tracce archeologiche.

Questa ipotesi non è falsificabile in senso stretto con i dati attuali: non si può dimostrare che qualcosa non sia mai esistito. Ma è corretto tenerla sul tavolo accanto all'ipotesi della tafonomia, perché le due hanno implicazioni di ricerca diverse. Se si pensa che l'arte c'era e sia sparita, ha senso investire in tecniche di rilevamento avanzate. Se si pensa che non ci fosse, ha più senso orientare la ricerca verso le fasi cronologiche successive, Neolitico ed età del Bronzo, dove la tradizione di arte rupestre all'aperto è attestata anche in contesti alpini comparabili. Il capitolo~\ref{cap:storia_territorio} di questo libro discute entrambe le fasi in prospettiva storica.
